\chapter{От эмбриона к золотому миру: \texttt{exp2} и \texttt{exp3}}

\section{\texttt{exp2}: эмбриональный граф из реальных источников}

\subsection{Смысл эксперимента}

\texttt{exp2} был первым настоящим столкновением с вопросом: можно ли вообще
собрать из реальных источников такой граф, который будет не просто лежать на
диске, а поддерживать полный Stage-1 pipeline. Речь шла не о красоте retrieval,
а о минимальной жизнеспособности всей линии:

\begin{enumerate}[leftmargin=2em]
\item seeded loading подграфа из реальных данных;
\item anchoring пользовательского запроса в этом графе;
\item диффузия активации;
\item упаковка результата в ограниченный контекст.
\end{enumerate}

Важное отличие от поздних экспериментов состоит в том, что здесь еще нет
полноценной ground-truth оценки качества ответа. Эксперимент отвечает на более
ранний вопрос: работает ли вообще ``физика системы'' на настоящем материале.

\subsection{Стенд}

Источники были следующими:

\begin{itemize}[leftmargin=2em]
\item полный дамп ConceptNet;
\item агрегированный ATOMIC 2019;
\item английский NumberBatch.
\end{itemize}

Загрузка строилась как seeded BFS от ограниченного набора входов с радиусом 3--4
и потолком числа concept-узлов. Дополнительно использовались:

\begin{itemize}[leftmargin=2em]
\item отдельный более жесткий порог для \texttt{RelatedTo};
\item точное лемматическое связывание ATOMIC с отключенным fuzzy fallback;
\item общий векторный словарь NumberBatch для якорения и similarity.
\end{itemize}

В терминах будущей архитектуры этот опыт проверял, может ли эмбрион
\emph{существовать как вычислимый объект}, а не лишь как схема в manual.

\subsection{Что здесь могло получиться}

На старте было две альтернативы.

\begin{enumerate}[leftmargin=2em]
\item \textbf{Оптимистическая.} Пайплайн в целом работает, а проблемы будут
касаться главным образом качества retrieval.
\item \textbf{Пессимистическая.} Реальные данные выявят системные дефекты еще до
момента, когда вообще можно говорить о retrieval-качестве: ошибки разметки,
схлопывание подписей, фантомные хабы, неработающие embedding lookup и т.\,п.
\end{enumerate}

\subsection{Что реально произошло}

Произошел более сильный и более полезный вариант пессимистической развилки:
pipeline в целом заработал, но именно на нем всплыл целый класс ошибок, которые
синтетический мир не мог бы показать.

\subsubsection*{(1) Ошибка labeling из POS-tagged URI}

Часть узлов получала содержимое вроде одной буквы \texttt{n}, потому что
последний сегмент URI принимался за label узла, хотя это был POS-tag. Ошибка
оказалась не косметической: в одном из запусков она затронула около трети узлов.

\subsubsection*{(2) Дублирование ребер}

Двунаправленное индексирование adjacency делало одно и то же ребро доступным из
обоих концов, а BFS потом повторно добавлял его в граф. Нужна была дедупликация
по assertion id самого дампа.

\subsubsection*{(3) Производство хабов связью \texttt{RelatedTo}}

Это одна из ключевых находок всей ранней серии. Оказалось, что проблема не в
``больших степенях вообще'', а в том, что именно \texttt{RelatedTo} на реальных
данных образует слишком широкие и семантически рыхлые маршруты. Отсюда уже в
\texttt{exp2} родилась мысль, что тип отношения в графе не является декоративной
меткой: он меняет реальную физику retrieval.

\subsubsection*{(4) Несовпадение формата NumberBatch и ConceptNet-URI}

English-only NumberBatch оказался записан как bare terms, а не как полные
ConceptNet URI. Поэтому все lookup молча проваливались, и embeddings не
подгружались вовсе.

\subsubsection*{(5) Переизбыточное ATOMIC-linking}

Даже после исправления NumberBatch слишком мягкая схема embedding-based linking
присоединяла к графу лавину эпизодов. Переход к exact lemma match дал резкое
сжатие числа episode nodes и тем самым показал, что на этом уровне важно не
``чуть лучше подобрать threshold'', а сначала вообще вырваться из режима
неконтролируемого фан-аута.

\subsection{Инфраструктурный результат}

После исправлений pipeline заработал end-to-end на реальных данных. Запрос
``what is the difference between lisbon and bacalhau'' успешно якорился,
распространял активацию и упаковывался в контекст.

Но этот результат сам по себе не означал, что retrieval хорош. Он означал более
скромную и при этом чрезвычайно важную вещь: архитектура вышла из стадии чистой
спекуляции и впервые столкнулась с реальной семантической грязью.

\begin{takeaway}
\texttt{exp2} доказал не качество памяти, а жизнеспособность эмбриона как
вычислимого объекта. Еще важнее он доказал, что реальные данные порождают
собственный класс багов, которых не существует на игрушечных графах. Это и
обосновывает всю дальнейшую линию: без реального эмбриона часть архитектурных
дефектов просто оставалась бы невидимой.
\end{takeaway}

\subsection{Что эксперимент оставил открытым}

Главный открытый вопрос после \texttt{exp2}: что значит ``хороший'' retrieval на
реальном графе, если там нет ground truth. Сомнительный ranking можно заметить,
но нельзя строго назвать неверным без контролируемой постановки.

Именно отсюда рождается \texttt{exp3}. Нужен был не новый retrieval law, а
маленький золотой мир, где правильный ответ заранее известен.

\section{\texttt{exp3}: золотой микромир как измеритель retrieval}

\subsection{Зачем понадобился второй стенд}

После \texttt{exp2} стало понятно, что два вопроса нельзя больше смешивать:

\begin{enumerate}[leftmargin=2em]
\item может ли pipeline работать на реальных данных;
\item выдает ли он правильный ответ, когда правильность вообще можно проверить.
\end{enumerate}

\texttt{exp3} не вводит новый retrieval-механизм. Он вводит новую
\emph{измерительную среду}: маленький hand-authored граф, где каждая неудача
может быть названа по имени.

\subsection{Конструкция золотого мира}

Граф был намеренно очень маленьким и семантически прозрачным. В нем были два
сюжетных кластера.

\begin{itemize}[leftmargin=2em]
\item Кофейный: \texttt{coffee}, \texttt{caffeine}, \texttt{beverage},
\texttt{hot drink}, \texttt{tea}, \texttt{juice}, \texttt{soda}.
\item Географический: \texttt{north dakota}, \texttt{south dakota},
\texttt{north}.
\end{itemize}

Смысл кластера про кофе был особенно важен. Он позволял сравнить:

\begin{itemize}[leftmargin=2em]
\item специфический, содержательно правильный ответ \texttt{caffeine};
\item широкий, структурно связный, но концептуально грубый distractor
\texttt{beverage}.
\end{itemize}

\subsection{Какие именно сценарии проверялись}

В \texttt{exp3} было четыре сценария.

\begin{enumerate}[leftmargin=2em]
\item \textbf{Multi-word anchoring.} Проверка, что \texttt{north dakota}
якорится как единый термин, а не распадается на \texttt{north} и
\texttt{dakota}.
\item \textbf{Distractor vs answer.} Ключевой сценарий:
``why coffee gives energy'' не произносит \texttt{caffeine} явно, поэтому
retrieval обязан выйти к нему за счет графовой структуры и similarity, а не за
счет банального подглядывания.
\item \textbf{Multi-anchor query.} Проверка, что разные темы могут жить вместе в
одном контексте.
\item \textbf{Regression on irrelevant words.} Проверка, что stopwords и
служебные слова не ломают anchoring.
\end{enumerate}

\subsection{Главный результат}

Сценарии 1, 3 и 4 прошли хорошо. Но второй сценарий дал первый по-настоящему
диагностический удар по retrieval-механике.

Для запроса про кофе и энергию peak activation выглядела так:

\begin{center}
\begin{tabular}{lcc}
\toprule
Узел & Амплитуда & Роль \\
\midrule
\texttt{beverage} & 146.67 & distractor \\
\texttt{coffee} & 100.00 & anchor \\
\texttt{hot drink} & 100.00 & anchor \\
\texttt{tea} & 100.00 & anchor \\
\texttt{caffeine} & 26.67 & правильный ответ \\
\bottomrule
\end{tabular}
\end{center}

Это не ``слегка неудачный ranking'', а качественно неправильное поведение:
широкий хаб побеждает специфический ответ более чем впятеро.

\subsection{Что именно это доказало}

Важен не только сам факт ошибки, но и ее структура. \texttt{caffeine} проиграл
не потому, что он был плохо представлен как embedding. Он проиграл потому, что
в базовой диффузии нет механизма, который бы делал различие между:

\begin{itemize}[leftmargin=2em]
\item узлом, в который сходится много путей просто по топологии;
\item узлом, который является более точным ответом на вопрос.
\end{itemize}

Это сразу породило рабочую гипотезу для \texttt{exp4}: возможно, хватает просто
добавить query-conditioned semantic gate, то есть умножать поток на релевантность
узла запросу.

\begin{riskbox}
Именно здесь серия могла бы свернуть в ложную сторону и начать считать, что
проблема целиком в embeddings. Но \texttt{exp3} устроен так, что дефект уже виден
на очень маленьком и прозрачном графе. Поэтому рабочая гипотеза после него --
не ``надо срочно менять векторное пространство'', а ``в operator diffusion
отсутствует релевантностный множитель''.
\end{riskbox}

\subsection{Связь \texttt{exp2} и \texttt{exp3}}

Пара \texttt{exp2}--\texttt{exp3} образует очень важное методологическое кольцо.

\begin{itemize}[leftmargin=2em]
\item \texttt{exp2} показал, что без реального контента не увидеть целый класс
инфраструктурных дефектов.
\item \texttt{exp3} показал, что без золотого мира нельзя всерьез говорить о
качестве retrieval, даже если pipeline уже технически работает.
\end{itemize}

Вместе они вводят основной стиль всей дальнейшей серии: чередование между
реализмом и контролируемостью.
